El desarrollo de Thary se apoyó mediante el uso de la herramienta de inteligencia artificial Claude (Anthropic), incluida su versión para el entorno de programación (Claude Code), como asistente de apoyo en diferentes etapas. El alcance de dicho uso se detalla a continuación:

\begin{itemize}
    \item Apoyo en la corrección ortográfica y gramatical de los capítulos, ofreciendo sugerencias de cómo reorganizar las secciones de la documentacion y redactando propuestas de texto a partir de las ideas y explicaciones del autor.
    \item Verificación constante del código fuente real del sistema Thary, con el fin de asegurar que las explicaciones y la documentación del documento coincidieran con las últimas versiones del código. Ya que debido a las constantes actualizaciones que recibió el sistema, fue común que varios capítulos necesitaran ser actualizados y reescritos para reflejar el estado actual del código.
    \item Apoyo en la escritura y creación de código de algunas funcionalidades del sistema, así como en la corrección rápida de errores de sintaxis y de ejecución.
    \item Apoyo en la corrección de errores de compilación y sintaxis del documento en \LaTeX.
\end{itemize}

Aun asi, se verificaron manualmente todas las aportaciones que dio la IA con el fin de asegurar que fuesen correctas, tambien se reviso el contenido generado para asegurar que no se introdujeran conceptos nuevos que no estan implementados en el sistema.